\documentclass[11pt]{article}
\usepackage[T1]{fontenc}
\usepackage{textcomp}
\usepackage[margin=0.75in]{geometry}
\usepackage{amsmath,amssymb,amsthm}
\usepackage{array,booktabs}
\usepackage{hyperref}
\setlength{\parindent}{0pt}
\newtheorem{theorem}{Theorem}
\theoremstyle{definition}
\newtheorem{definition}{Definition}
\title{Flow Proof Artifact: Ring-one-plus-one-derived}
\author{Generated by \texttt{flow doc proof}}
\date{Flow Proof Book}
\begin{document}
\maketitle
One plus one is two in a ring witness.\\[0.5em]
\textbf{Source.} dummit-foote — *Abstract Algebra*, §7.1\\[0.5em]
\bigskip
\noindent{\large \textbf{Derived fact 1.} \textit{1 + 1 = 2} \hfill Ring \cdot addition \cdot one plus one derived}\par
\smallskip\noindent\textit{Coordinate: Ring \cdot addition \cdot one plus one derived}\par
\smallskip\noindent\textit{Source: dummit-foote}\par
\smallskip\noindent\textit{Built on: one plus one is two, for addition on Ring}\par
\medskip
\noindent\textbf{Goal.} 1 + 1 = 2\par
\noindent$\displaystyle 1 + 1 = 2$\par
\medskip
\noindent
\begin{tabular}{@{} >{\raggedright\arraybackslash}p{0.06\textwidth} >{\raggedright\arraybackslash}p{0.47\textwidth} >{\raggedleft\arraybackslash}p{0.41\textwidth} @{}}
\textbf{\#} & \textbf{Proof} & \textbf{Mathematics} \\
\midrule
\textbf{1.} & We prove that one plus one derived for addition on Ring. & \\[0.45em]
\textbf{2.} & We invoke the derived fact governing addition on Ring: one plus one is two, for addition on Ring. & \\[0.45em]
\textbf{3.} & From step 2, this implies 1 plus 1 equals 2. Hence proven. & $\displaystyle 1 + 1 = 2$ \\[0.45em]
\bottomrule
\end{tabular}
\label{thm:Ring-addition-one_plus_one_derived}
\par\medskip\hrule
\end{document}
